\clearpage\thispagestyle{empty}
\pdfbookmark[0]{Quellen zu Juno (23. Dezember 1813)}{juno1813dec23-sources}
\begin{center}\textbf{Quellen zu Juno (23. Dezember 1813)}\end{center}
\noindent Diese KI-Übertragung ins heutige Deutsch folgt dem Bericht mit der Überschrift Juno in den \href{https://gdz.sub.uni-goettingen.de/id/PPN236007467}{\textit{Werken}, Band~VI} von Gauß (Göttingen, 1874), S.~372--373. Er erschien zuerst unter Göttingen in den \textit{Göttingischen gelehrten Anzeigen}, Stück~204 (23. Dezember 1813), S.~\href{https://api.digitale-sammlungen.de/iiif/image/v2/bsb10538684_01009/full/1600,/0/default.jpg}{2033}, \href{https://api.digitale-sammlungen.de/iiif/image/v2/bsb10538684_01010/full/1600,/0/default.jpg}{2034}. Die Links führen zum digitalisierten Exemplar der Bayerischen Staatsbibliothek.
\par\smallskip\noindent Im Satz zur Berechnung der Opposition steht in der Werkausgabe der Plural, im Erstdruck der Singular. Beide geben für die Oppositionszeit 5 Sekunden an und setzen am Ende der Beobachtungszeit ein Doppelstrichzeichen. In der gesonderten Dezemberkorrespondenz beträgt die Sekundenangabe zur Opposition 25. Die \textit{Georgsharfe} ist auf \href{https://www.e-rara.ch/i3f/v20/3341782/full/2200,/0/default.jpg}{Tafel XII} von Johann Elert Bodes \textit{Uranographia} (Berlin, 1801) als \textit{Harpa Georgii} eingezeichnet, ETH-Bibliothek Zürich, Rar 10278. KI-Übertragung und KI-Satz.
\clearpage
